\section{Conclusion and perspectives}
\label{sec:conclusion}

We have derived an exact partial-sum calculation for triangular reaction--diffusion systems with critical gradient growth. The exact tested identity contains the full gradients of preceding components. If these gradients are uniformly bounded in $L^2(Q_T)$, Young's inequality with a free parameter absorbs all cross terms for any fixed positive diffusion vector; no diffusion ordering is intrinsic to this conditional step. Conversely, bounds on $\nabla T_k(u_{j,n})$ alone do not justify the same conclusion. The main analytical contribution is the identification of this exact full-gradient cross-term structure and the resulting conditional coercivity estimate for any fixed positive diffusion vector. An important consequence is the explicit separation between the algebraic absorption step and the compactness theory that must still be supplied by the chosen solution framework.

The numerical contribution now has three complementary levels. First, an independent manufactured-solution experiment verifies the implementations of the classical solvers: spatial refinement gives approximately second-order RMSE behaviour for both FDM and the assembled $P_1$ FEM, while separate temporal refinement gives first-order behaviour for the semi-implicit update. Second, the two-component benchmark retains the reproducible PINN--FDM--FEM comparison, common-point independent residual testing, five-seed descriptive statistics, an explicit $4.243\%\pm0.378\%$ relative initial-data mismatch, mass and energy diagnostics, and collocation ablation. On that smooth forward problem, the classical methods remain substantially more accurate and cheaper to construct than the PINN, while the PINN provides a nonnegative boundary-compatible differentiable surrogate. Third, the new three-component benchmark uses the non-monotone diffusion vector $(0.1,0.8,0.3)$ and realizes two simultaneous cross terms in a controlled structural test. The preceding-gradient bounds are established independently as $C_1\le1.2$ and $C_2\le0.0708984375$, yielding a residual Young contribution bounded by $0.321357421875$. A post-processed coercivity-bound consistency diagnostic over four truncation levels gives $\mathcal C_k/\mathcal B_k<0.221$ for the refined reference and essentially the same ratios for FDM and FEM. Against the refined numerical reference, the three-component FDM and FEM attain RMSE values $1.48\times10^{-4}$ and $1.14\times10^{-4}$, respectively, while remaining nonnegative and exhibiting non-increasing partial masses on the common evaluation grid.

The scope is therefore precise. Proposition~\ref{prop:partial-sum} is a conditional coercivity result whose full-gradient hypothesis must come from the surrounding analytical framework; the three-component experiment provides a controlled one-dimensional realization of that structure. The numerical comparison is correspondingly diagnostic: it separates field error, residual generalization, initial-data mismatch, and structural observables on the tested forward problem. Extension to a complete existence theorem, higher-dimensional or inverse settings, or a calibrated application requires the additional analytical and modelling ingredients identified below.

\subsection*{Perspectives}

Three directions follow naturally. On the analytical side, the priority is to derive the full preceding-gradient estimates from independently checkable hypotheses, or to replace them by a duality, entropy, or renormalized argument that closes the compactness proof. Weighted lower-triangular structures, no-flux boundary conditions, uniqueness, and long-time behaviour are further relevant extensions. On the numerical side, FP64 Adam--L-BFGS, loss-balancing strategies motivated by known gradient-flow pathologies~\cite{wang2021gradient,wang2022ntk}, causality-aware training~\cite{wang2024causality}, and adaptive residual sampling~\cite{wu2023sampling} should be compared under matched computational budgets; two-dimensional domains, stiff diffusion contrasts, inverse identification, and residual-based a posteriori indicators would test whether the differentiable surrogate offers a genuine advantage. On the modelling side, the transfer--dissipation chain can be specialized to a selected interfacial chemical, biological, materials, or image-processing problem only after dimensional consistency and measurable observables have been specified. These extensions would turn the present auditable benchmark into either a complete analytical theorem or a validated application study.

\subsection*{Code and data availability}
The supplementary reproducibility archive contains the code, fixed configurations and seeds, raw metrics, solution arrays, independent-residual, initial-condition and ablation outputs, refinement data, dependency specification, and figure-generation routines used for the reported results. The archived scripts regenerate the validation figures and post-processed diagnostics from the stored outputs without requiring the solvers to be rerun.

\subsection*{Acknowledgments}
The authors acknowledge the computational resources provided by the National School of Applied Sciences of Safi, Cadi Ayyad University, Morocco.